% Proposed after checkpoint09; actual quotient frames, not a full scalar lift.
\statement{A.66}{Proved}{A class-three twenty-direction quotient admits compatible involutory peripheral frames}
After another permitted choice of the discretionary affine data, the
integral cross-cup tower of A.65 has an actual further source quotient
\[
 1\longrightarrow (V^{(3)})^{20}\longrightarrow E^{(3)}
 \longrightarrow E\longrightarrow1.
 \tag{A.66.1}
\]
Here \(V^{(3)}\) has the cubed monomial scalars. The twenty coordinate
classes are k-linearly independent and nonsplit. The actual quotient
admits a positive peripheral automorphism, an involutory negative
companion, and involutory s- and r-normalizer frames retaining the
construction-selected initial C-data and older families. The new
positive action is the identity on its entire free-cusp group factor.
The scalar refinement above (A.66.1), and realization of these exact
new frames by a full automorphism, are not asserted here.

\begin{proof}
\emph{An actual class-three group and source map.}\par
Put \(A=V^4\), \(T=(V^{(2)})^6\),
\(U=(V^{(3)})^4\), and \(Z=(V^{(3)})^{16}\).
The indices of T are \(i<j\), those of U are \(i<j<k\), and
those of Z are ordered pairs \((r,i)\), with all indices from one
to four. On \((a,t,u,z)\) define multiplication by
\[
 \begin{aligned}
 a''&=a+b,\\
 t''_{ij}&=t_{ij}+t'_{ij}+a_i*b_j,\\
 u''_{ijk}&=u_{ijk}+u'_{ijk}+a_i*t'_{jk}+t_{ij}*b_k,\\
 z''_{ri}&=z_{ri}+z'_{ri}+a_r^{[2]}*b_i.
 \end{aligned}
 \tag{A.66.2}
\]
The product \(*\) and the power in brackets are coordinatewise.
The two triple-cup coboundaries in the u formula cancel. The other
corrections are biadditive over \(\mathbb F_2\), so (A.66.2) is
associative. H acts strictly with scalar degrees one, two, three,
and three on the four displayed blocks. Denote this finite two-group
by K. Its third lower-central subgroup is U. Its derived and
Frattini subgroups are \(T\oplus U\oplus Z\), and its Frattini
quotient is A. To check attainment, single-row brackets give T,
the triples \([a_i,[a_j,a_k]]\) give each U-coordinate, and the
ordinary brackets give
\[
 a_r^{[2]}*b_i+b_r^{[2]}*a_i.
 \tag{A.66.3}
\]
The functions \(a^2b\) and \(ab^2\) are independent as
\(\mathbb F_2\)-bilinear functions on \(k^2\), and
\(a^2b+ab^2\ne0\). Thus all sixteen ordered Z-coordinates
are attained independently. Their cross-Frobenius support
distinguishes them from the old same-family T brackets.

Use the four normalized cusp cocycles, whose s-value is zero.
In \(K\rtimes C\), take S pure and R conjugate to pure r by
an element with abelian coordinate \(r e(r)\) and all higher
coordinates zero. This gives \(S^2=R^3=1\), and it projects
literally onto the source map of A.63. Initial-base
\(H^1(C,V)=0\) makes the actual R-image onto A. The finite
two-group Frattini theorem makes it all K.

Pairing with the full old integral residual over A gives the
actual fiber product: its Frattini quotient is the complete
\(V^8\), onto which the actual R-image already surjects.
This proves (A.66.1), with no presumed finite-quotient realization.
The Frobenius weight identity \(3\mathcal S=-4\mathcal S\)
identifies the underlying binary simple type with type 108;
we nevertheless keep the literal cubic coefficient action.

For independence, push out by a nonzero k-linear functional on
the twenty coordinates. If its U-part is nonzero, the resulting
residual still has nonzero third commutators. A section from the
old class-two residual would be onto modulo Frattini and hence
onto, an impossibility. If its U-part is zero, use the ordinary
brackets (A.66.3). An H-map from old T to \(V^{(3)}\) sends a
vector supported at one physical weight to the opposite weight,
with its entries squared; this follows from the three disjoint
Frobenius weight families. An H-map from the old trivial center
is zero. On two inputs at a single physical weight, the proposed
ordinary bracket stays at that weight, so it cannot factor through
the old derived group. Independence of \(a^2b\) and \(ab^2\)
then forces every ordinary coefficient to vanish. A splitting induces an
H-equivariant map on the old residual derived group: the new
kernel centralizes the residual, so changing the chosen lifts
of C cannot alter that derived action. This proves the
twenty independent nonsplit directions.

\emph{Exact action lifts.}\par
For a six-by-four coefficient matrix L, the old shear
\(t_{ij}\mapsto t_{ij}+\sum_rL_{ij,r}a_r^{[2]}\) lifts to the
following H-equivariant automorphism of K: fix a and z and set
\[
 u_{ijk}\longmapsto u_{ijk}
   +\sum_r\left\{L_{ij,r}z_{rk}
        +L_{jk,r}(z_{ri}+a_i*a_r^{[2]})\right\}.
 \tag{A.66.4}
\]
Indeed, the coboundary of \(z_{rk}\) is \(a_r^{[2]}*b_k\),
whereas that of \(z_{ri}+a_i*a_r^{[2]}\) is
\(a_i*b_r^{[2]}\). These are precisely the change in the
u product after the t shear. The resulting maps \(F_L\)
satisfy \(F_LF_{L'}=F_{L+L'}\).

The lower canonical-c correction has only four cusp cohomology
coefficients, together with an inner term. Formula (A.66.4),
the actual monomial c-action, and these inner terms therefore
give an honest initial lift. Correct it by the complete source
one-cocycle space in coefficient \(V^{(3)}\). The intertwiner
from V is
\[
 j_3(v)_\lambda=(v_{-\lambda})^4.
 \tag{A.66.5}
\]
It gives the same 44 independent source-value columns as before.
Matching \(c(s)=s\) and \(c(r)=sr^{-1}s\) reconstructs the
literal canonical source action on the actual quotient. Its
square is one because these are the source generators.

\emph{A compatible old affine choice.}\par
Impose \(W_{\rm cusp}=0\), retaining the old central class
equations but allowing the free s-frame t to vary. The exact
system has rank 102 in 360 binary coordinates. On its
258-dimensional homogeneous space, the old central norm has
polar radical 231, radical image zero, and a 27-bit quadratic
core. The exact solution numbered 2128959 was found and checked
in the original group law. Its initial local frame was then
adjusted within that radical. The two opposite-weight local gates
have 29 possibilities each; 29 of their 841 paired possibilities
are compatible with this same norm core. The radical has local
image rank 24. The saved choice has local indices \([6,0]\).

These are finite existence certificates, not dimension heuristics.
The original-law verified tuple is
\path{work_A/A66_W0_local_compatible_partial.json}.
The relevant producers and exact coefficient records are
\path{work_A/check_A66_zero_W_free_t_norm.py},
\path{work_A/A66_zero_W_free_t_quadratic_core.json},
\path{work_A/check_A66_W0_radical_local_gate.py}, and
\path{work_A/check_A66_W0_local_adjustment.py}.
The lower-six and ordinary-sixteen direct group checks are
\path{work_A/check_A66_W0_local_lower.out} and
\path{work_A/check_A66_W0_local_ordinary.out}.
They verify their negative squares and complete s-frame norms.
Their positive correction coefficients are identically zero.

Consequently the positive map on K is literally the identity,
paired with the retained old action over C. This is an actual
automorphism because the old action fixes the finite H-head and
the four cusp cocycles. Its negative companion on the K-factor
is the canonical c just constructed. Its square on the whole
fiber product is therefore one. This does not identify a full
source lift of this selected new automorphism.

\emph{The remaining four frame norms.}\par
Combine the saved lower T-frame and ordinary Z-frame, initially
putting the four U-frame coordinates equal to zero. It already
has the correct s-conjugacy, and its norm lies in U. That norm
is fixed by both S and the coefficient involution D.

On the new kernel, canonical c has upper-triangular multiplicity
map \(I+J_{L_c}\), followed by its cubic monomial action.
The map \(J_{L_c}:k^{16}\to k^4\) has rank four. The exact
certificate gives a triangular four-by-four minor with diagonal
\((7,2,2,2)\) in the fixed F8 encoding. A homogeneous ordinary
correction Z fixed by S and D has norm
\[
 (J_{L_c}Z,0)\in U\oplus Z.
 \tag{A.66.6}
\]
A constant right inverse of \(J_{L_c}\) commutes with the common
monomial actions. It therefore lifts every joint-invariant U
residue to a joint-invariant ordinary correction. This kills all
four remaining norms while preserving the previous lower and
ordinary equations. It changes the ordinary frame only within
its homogeneous norm-zero family.

The complete direct witness is
\path{work_A/A66_W0_twenty_sframe_solution.json}; its producer
and stdout are
\path{work_A/check_A66_W0_twenty_sframe.py} and
\path{work_A/check_A66_W0_twenty_sframe.out}.
The independent coefficient check is
\path{work_CE/code/check_A66_UT4_frame_surjection.py}.
These verify source relations, canonical c, literal cusp,
negative square, s-conjugacy, and the exact s-normalizer norm.

\emph{The r-frame and the retained initial choice.}\par
We make explicit an additional closed condition available in
the A.33 construction. Its dyadic s-frame comes from the original
full Galois involution. Its dyadic r-reflection is still purely
conjugate to that involution's r-reflection: A.16 retains the
chosen reflection, and the final coprime correction conjugates it.
For the common full Galois involution, the two reflections
\(s d_s\) and \(d_r\) are purely conjugate. For example, if
\(b=\nu c\nu^{-1}\), with
\(\nu(s)=u_s s u_s^{-1}\) and
\(\nu(r)=u_r r^{\epsilon}u_r^{-1}\), their conjugator is
\(u_r^{-1}u_s\). Composing gives the dyadic transporter used here.

At each finite congruence level retain its finite transporter
image. At every odd relative prime, two involutions with that
same image are complements to a pro-odd kernel, hence are
conjugate by that kernel. The transporter therefore glues across
the relative fiber product. Add this transporter as a compact
variable in A.33. For every finite collection use a common level;
the finite-intersection property is unchanged. The initial
C-frames can thus be selected with a pure transporter from
\(s_Cd_{s,C}\) to \(d_{r,C}\). This is a selection from the
actual construction, not a claim about arbitrary norm-zero frames.

First use the new actual quotient over C before adjoining the
older semisimple families; its kernel is pro-two. The solved
s-reflection commutes with s, so \(s d_s\) is an involution.
Lift the selected C-transporter. Its conjugate is an involution
over \(d_{r,C}\). Apply A.16 in the preimage of the actual
finite subgroup \(\langle r_C,d_{r,C}\rangle\simeq S_3\),
retaining this involution. Coprime conjugacy of its order-three
complement onto the actual \(\langle r\rangle\) supplies an
involutory r-normalizer over the exact original C-frame.
Pair this r-frame with the retained older r-frames over their
identical C-frame, using the already proved actual fiber product
and separation of residual simple types. No pro-two assertion
is made about the whole joint quotient with those older families.
This is an existence proof; no coordinates of that r-frame are
asserted in the displayed computational witness.

Finally, construct the integral actions, rather than infer them
from a particular full lift. Write \(P_n^{\rm free}\) for the
four-free-cusp class-two group. The actual source comparison and
the exact centers of A.58--A.62 give
\[
 P_n\simeq P_n^{\rm free}
       \mathbin{\times}_{P_1^{\rm free}}P_1.
 \tag{A.66.7}
\]
Only the six free off-diagonal center coordinates acquire higher
two-power exponents; all mixed and binary center terms are
already modulo two. Here \(P_n\) denotes the residual group.
The same comparison of source maps extends (A.66.7) to the
corresponding actual quotients over C:
\[
 D_n\simeq D_n^{\rm free}
       \mathbin{\times}_{D_1^{\rm free}}D_1.
\]
The source C4 carry is retained in the full \(D_1\)-factor.
In \(D_n^{\rm free}\), S is pure and the chosen
opposite-pair orientation is S-invariant. These observations
verify the actual source maps; alternatively the surjective
residual map on the complete abelianization and exact center
gives the same identification.

The chosen old action is literally identity on
\(P_1^{\rm free}\). Its free abelian action is identity and
its free inner vector is zero; the remaining possible central
character vanishes on s by the strict S-frame conjugacy and
on r by order three. Pair identity on \(P_n^{\rm free}\)
with the chosen action on \(P_1\). This gives compatible
actual actions with zero free inner vector at every n.
The invariant-lifting and orientation argument of A.62 lifts
the nonzero free s-frame through those levels; it does not
deduce an integral norm from a mod-two norm. Old scalar-half
corrections have no K-coordinate. Finite common refinements
and compactness retain all levels and older families; the
latter still have different residual abelian simple types.
All statements concern the actual quotient and its frames. The
next scalar lifting criterion remains to be checked.
\end{proof}
